% Claim proof file for row claim_3_16 -- "MarginalizationPreserves".
% Auto-generated stub by scaffold/solving_current_row.py at row start. A
% manager agent dedicated to writing the TeX proof fills in the body below.
% The proof must:
%   - Stay close to the lecture notes' proof if one exists (always search
%     the LN first for a `\begin{proof}` block immediately following the
%     claim; copy / lightly edit when present).
%   - Otherwise use the LN paradigm and cite prior claims by their `ref`.
%   - Be self-contained enough that a mathematician can follow it without
%     re-reading the LN.
%   - Have no `sorry`, `omitted`, or "obvious" shortcuts.
%
% Compiles standalone via the `subfiles` package.

\documentclass[main]{subfiles}

\begin{document}
% ref: claim_3_16 (proof, with statement restated for readability)
% title: MarginalizationPreserves
\def\rowref{claim\_3\_16}
\phantomsection\label{claim_3_16}
%
% Cross-references: see claim_<ref>_statement_<title>.tex in this folder
% for the corresponding statement.

% --- statement (restated) ---------------------------------------------------
\begin{Rem}[MarginalizationPreserves]\label{rem:marg_preserves_ancestors_bifurcations_acyclicity}
Marginalization preserves ancestral relations, bifurcations and acyclicity:
    \begin{enumerate}
        \item For $v_1,v_2 \in G$ with $v_1,v_2 \notin W$ we have the equivalence:
            \[ v_1 \in \Anc^G(v_2)\quad \iff \quad v_1 \in \Anc^{G^{\sm W}}(v_2).   \]
          \item\label{rem:marg_preserves_bifurcations} 
            For $v_1,v_2 \in G \sm W$ (and, optionally, $v_3 \in G\sm W$): there is a bifurcation between $v_1$ and $v_2$ (with source $v_3$) in $G$ if and only if there is a bifurcation between $v_1$ and $v_2$ (with source $v_3$) in $G^{\sm W}$.
        \item If the CDMG $G$ is acyclic then so is $G^{\sm W}$ and a topological order of $G$ induces a
            topological order on $G^{\sm W}$ (by just ignoring the nodes from $W$).
    \end{enumerate}
\end{Rem}

% --- proof ------------------------------------------------------------------
% The LN supplies a proof at lecture-notes/lecture_notes/graphs.tex:976--993
% that covers only item 2 (bifurcations), and does so by induction on |W| via
% \refrow{claim_3_17} (Marginalizations Commute) -- a forward reference. The
% proof below replaces that induction by a direct multi-vertex shrink/expand
% argument through the simp-characterisations of \refrow{def_3_14}, and adds
% from-scratch proofs of items 1 (ancestors), 3a (acyclicity) and 3b
% (topological order); the LN leaves items 1 and 3 implicit.
\begin{proof}
\textit{Setup -- walk translators.}  Fix a CDMG $G = \lt J, V, E, L \rt$ and a set $W$ of vertices; write $G^{\sm W}$ for the marginalization \refrow{def_3_14}. The directed and bidirected edges of $G^{\sm W}$ are characterised by the corresponding simp-lemmas of \refrow{def_3_14}: a pair $(u, v)$ is in $E^{\sm W}$ iff $u \in J \cup (V \sm W)$, $v \in V \sm W$, and there exists a directed walk $u \tuh \cdots \tuh v$ in $G$ of length $\ge 1$ with every intermediate vertex in $W$; the pair $(u, v)$ is in $L^{\sm W}$ iff $u, v \in V \sm W$, $u \ne v$, neither $(u, v)$ nor $(v, u)$ is realised by a directed walk in $G$ with intermediate vertices in $W$, and some bifurcation between $u$ and $v$ in $G$ (in either walk direction) has all intermediate vertices in $W$. The two negated-directed-walk clauses are Lean-encoding artefacts forced by $E \cap L = \emptyset$ in \refrow{def_3_1}; the symmetric ``either walk direction'' in the bifurcation clause matches the LN's symmetric ``between $v_1, v_2$''.

The proof uses two constructions repeatedly:

\smallskip
\textbf{(S) Shrink.} Given a directed walk $\pi : u \tuh \cdots \tuh v$ in $G$ of length $\ge 1$ with $u, v \notin W$, list the vertices of $\pi$ \emph{not in $W$} in walk order as $u = w_0, w_1, \dots, w_r = v$; by construction all vertices strictly between $w_i$ and $w_{i+1}$ on $\pi$ lie in $W$. Each block $w_i \tuh \cdots \tuh w_{i+1}$ is a directed walk in $G$ of length $\ge 1$ with interior in $W$, so $(w_i, w_{i+1}) \in E^{\sm W}$. Concatenating these edges yields a directed walk $\pi^{\sm W} : u \tuh \cdots \tuh v$ in $G^{\sm W}$, of length $r \ge 1$. Crucially, this is a \emph{single-step} multi-vertex shrinkage -- no iteration over $|W|$ is required.

\smallskip
\textbf{(E) Expand.} Conversely, given a directed walk $\pi' : u \tuh \cdots \tuh v$ in $G^{\sm W}$, each step $(u_i, u_{i+1}) \in E^{\sm W}$ unfolds to a directed walk in $G$ from $u_i$ to $u_{i+1}$ of length $\ge 1$ with interior in $W$. Concatenating gives a directed walk $\pi'^{\uparrow} : u \tuh \cdots \tuh v$ in $G$, of length $\ge \pi'.\mathrm{length}$.

Both translators preserve endpoints and the all-$\tuh$ shape of a directed walk; reversing turns each into a translator for all-$\hut$ walks, so they apply to both arms of a bifurcation.

\subsubsection*{Item 1 -- ancestral relations.}
\textit{(Justifies \texttt{marginalize\_anc\_iff}.)} Let $v_1, v_2 \in G$ with $v_1, v_2 \notin W$. By \refrow{def_3_5}, $v_1 \in \Anc^H(v_2)$ for a CDMG $H$ iff $v_1 \in H$ and a directed walk $v_1 \tuh \cdots \tuh v_2$ exists in $H$. The hypothesis $v_1 \notin W$ gives $v_1 \in G$ iff $v_1 \in J \cup (V \sm W) = J^{\sm W} \cup V^{\sm W}$ iff $v_1 \in G^{\sm W}$, so the only content of the biconditional is the directed-walk existential.

$(\Rightarrow)$ Given a directed walk $\pi : v_1 \tuh \cdots \tuh v_2$ in $G$. If $\pi.\mathrm{length} = 0$, then $v_1 = v_2$ and the trivial $\mathrm{Walk}.\mathrm{nil}$ at $v_1$ in $G^{\sm W}$ witnesses the conclusion. Otherwise apply (S) to $\pi$.

$(\Leftarrow)$ Given a directed walk $\pi' : v_1 \tuh \cdots \tuh v_2$ in $G^{\sm W}$. If $\pi'.\mathrm{length} = 0$, take the trivial $\mathrm{Walk}.\mathrm{nil}$ at $v_1$ in $G$; else apply (E).

\subsubsection*{Item 2 -- bifurcations (no source).}
\textit{(Justifies \texttt{marginalize\_bifurcation\_iff}.)} The LN's parenthetical ``(with source $v_3$)'' is intentionally \emph{deferred}: the Lean-side $L^{\sm W}$ exclusion clause can transform a bidir-hinge bifurcation in $G$ (no source) into a backward-hinge bifurcation in $G^{\sm W}$ (with source) by absorbing the hinge into a directed shortcut. The symmetric ``between $v_1, v_2$'' reading absorbs that case cleanly for the no-source half (the absorbed bifurcation reappears in the opposite walk direction of the symmetric $\vee$); the with-source half does not, and is studied separately. See workspace risk section~5.2 attached to \refrow{claim_3_16}.

Fix $v_1, v_2 \in G \sm W$. Following the symmetric reading, we prove that
\[
   (\exists\,\pi : v_1 \rightsquigarrow v_2,\ \pi.\mathrm{IsBifurcation}) \ \vee\ (\exists\,\pi : v_2 \rightsquigarrow v_1,\ \pi.\mathrm{IsBifurcation}) \quad\text{in } G
\]
is equivalent to the same disjunction with $G$ replaced by $G^{\sm W}$. By symmetry on $v_1, v_2$, it suffices to convert any bifurcation
\[
   \pi : v_1 \hut u_1 \hut \cdots \hut u_{k-1} \hus u_k \tuh \cdots \tuh u_{n-1} \tuh v_2 \qquad\text{in } G
\]
into a bifurcation between $v_1$ and $v_2$ in $G^{\sm W}$ in one walk direction (and conversely). The hinge $u_{k-1} \hus u_k$ is either bidirected ($u_{k-1} \huh u_k$) or directed backward ($u_{k-1} \hut u_k$, apex at $u_k$); the reversed left arm $u_{k-1} \tuh \cdots \tuh v_1$ and the right arm $u_k \tuh \cdots \tuh v_2$ are directed walks in $G$.

\smallskip
\noindent$(\Rightarrow)$ Let $m$ be the closest vertex to $u_{k-1}$ on the directed walk $u_{k-1} \tuh \cdots \tuh v_1$ that lies \emph{outside} $W$, and $m'$ the closest such vertex to $u_k$ on $u_k \tuh \cdots \tuh v_2$; both exist because $v_1, v_2 \notin W$.

\textbf{Bidir-hinge case ($u_{k-1} \huh u_k$).} Apply (S) to the directed walks $m \tuh \cdots \tuh v_1$ and $m' \tuh \cdots \tuh v_2$ to obtain $\pi_L^{\sm W} : v_1 \hut \cdots \hut m$ and $\pi_R^{\sm W} : m' \tuh \cdots \tuh v_2$ in $G^{\sm W}$ (each possibly trivial if $m = v_1$ resp.\ $m' = v_2$). The sub-walk
\[
   q :\ m \hut \cdots \hut u_{k-1} \huh u_k \tuh \cdots \tuh m' \qquad\text{in } G
\]
is a bifurcation between $m$ and $m'$ with interior in $W$ (intermediate vertices are $u_{k-1}, u_k$ if in $W$, plus all the in-$W$ vertices on the two arms between $m, m'$ and the hinge), so by the $L^{\sm W}$ characterisation either $(m, m') \in L^{\sm W}$ -- in which case the splice
\[
   v_1 \hut \cdots \hut m \huh m' \tuh \cdots \tuh v_2 \qquad\text{in } G^{\sm W}
\]
is a bifurcation between $v_1$ and $v_2$ (first disjunct) -- or the exclusion fires, namely a directed walk through $W$ in $G$ realises $(m, m')$ or $(m', m)$ as a member of $E^{\sm W}$. If $(m, m') \in E^{\sm W}$, read the constructed walk in $G^{\sm W}$ in the reverse direction $v_2 \to v_1$:
\[
   v_2 \hut \cdots \hut m' \hut m \tuh \cdots \tuh v_1 \qquad\text{in } G^{\sm W},
\]
a bifurcation between $v_2$ and $v_1$ with backward hinge $m' \hut m$ (apex at $m$), landing in the second disjunct. The mirror case $(m', m) \in E^{\sm W}$ yields a first-disjunct bifurcation $v_1 \hut \cdots \hut m \hut m' \tuh \cdots \tuh v_2$ by the symmetric splice.

\textbf{Backward-hinge case ($u_{k-1} \hut u_k$, apex at $u_k$).} Both arms of $\pi$ extend to directed walks rooted at $u_k$: $u_k \tuh u_{k-1} \tuh \cdots \tuh v_1$ and $u_k \tuh \cdots \tuh v_2$. If $u_k \notin W$, apply (S) to each (rooted at $u_k$) and splice at $u_k$ to obtain the backward-hinge bifurcation $v_1 \hut \cdots \hut u_k \tuh \cdots \tuh v_2$ in $G^{\sm W}$, landing in the first disjunct.

If $u_k \in W$, the local sub-walk
\[
   q' :\ m \hut \cdots \hut u_{k-1} \hut u_k \tuh \cdots \tuh m' \qquad\text{in } G
\]
is \emph{structurally the same as the bidir-hinge case's local bifurcation $q$, just with a backward hinge $u_{k-1} \hut u_k$ at apex $u_k$ instead of a bidirected hinge}: it is a bifurcation between $m$ and $m'$ in $G$ with interior in $W$ (intermediate vertices: those strictly between $m$ and $u_{k-1}$ on the left, the hinge endpoint $u_{k-1}$ if $u_{k-1} \in W$, the apex $u_k$, and those strictly between $u_k$ and $m'$ on the right -- all in $W$ by definition of $m, m'$ and the sub-case hypothesis $u_k \in W$). The same three-way analysis as in the bidir-hinge case, applied to $q'$ through the $L^{\sm W}$ characterisation, then yields:
\begin{itemize}
\item $(m, m') \in L^{\sm W}$ (no exclusion fires): first disjunct, via the splice $v_1 \hut \cdots \hut m \huh m' \tuh \cdots \tuh v_2$ in $G^{\sm W}$ (a bidirected hinge in $G^{\sm W}$ in spite of the backward-hinge origin in $G$ -- which is precisely the no-source case the symmetric $\vee$ is designed to absorb);
\item exclusion fires as $(m, m') \in E^{\sm W}$: second disjunct, via the reverse read $v_2 \hut \cdots \hut m' \hut m \tuh \cdots \tuh v_1$ in $G^{\sm W}$, a backward-hinge bifurcation with apex at $m$;
\item exclusion fires as $(m', m) \in E^{\sm W}$: first disjunct, via $v_1 \hut \cdots \hut m \hut m' \tuh \cdots \tuh v_2$ in $G^{\sm W}$, a backward-hinge bifurcation with apex at $m'$.
\end{itemize}
The same exclusion-direction-flip phenomenon as in the bidir case is therefore at work; the symmetric $\vee$ is what lets the second-disjunct landing escape the no-with-source constraint cleanly.

\smallskip
\noindent$(\Leftarrow)$ Suppose, WLOG,
\[
   \pi' : v_1 \hut u_1 \hut \cdots \hut u_{k-1} \hus u_k \tuh \cdots \tuh u_{n-1} \tuh v_2 \qquad\text{in } G^{\sm W}.
\]
Apply (E) to the reversed left arm $u_{k-1} \tuh \cdots \tuh v_1$ and to the right arm $u_k \tuh \cdots \tuh v_2$ to obtain directed walks in $G$ with intermediate vertices in $W$. Expand the hinge by cases:
\begin{itemize}
\item Backward hinge $u_{k-1} \hut u_k \in E^{\sm W}$: by the $E^{\sm W}$ characterisation, a directed walk $u_k \tuh \cdots \tuh u_{k-1}$ in $G$ with interior in $W$ exists; splice it in.
\item Bidirected hinge $u_{k-1} \huh u_k \in L^{\sm W}$: by the $L^{\sm W}$ characterisation, a bifurcation
\[
   u_{k-1} \hut w_1 \hut \cdots \hut w_{k'-1} \hus w_{k'} \tuh \cdots \tuh w_{n'-1} \tuh u_k \qquad\text{in } G
\]
exists with all $w_j \in W$ (possibly none, in which case the bifurcation degenerates to the bidirected or backward edge of $G$ directly connecting $u_{k-1}, u_k$ -- both possibilities are admitted as a length-$1$ ``hinge bifurcation'' by \refrow{def_3_4}, which allows empty arms and a single hinge edge; \refrow{def_3_14}'s $\vee$-of-two-walk-directions clause supplies either walk direction for the witness). Splice this in place of the hinge.
\end{itemize}
The all-$\hut$ orientation of the left arm and the all-$\tuh$ orientation of the right arm survive expansion (each $G^{\sm W}$-step expands to a directed walk of the matching orientation), and the spliced-in hinge contributes a $\hut\cdots\hut\hus\tuh\cdots\tuh$ V-shape lying entirely inside the original hinge slot. The composite is therefore a bifurcation between $v_1$ and $v_2$ in $G$.

\subsubsection*{Item 3a -- acyclicity.}
\textit{(Justifies \texttt{marginalize\_isAcyclic}.)} Assume $G$ is acyclic (\refrow{def_3_6}). Suppose for contradiction $v \in G^{\sm W}$ and $\pi' : v \tuh \cdots \tuh v$ in $G^{\sm W}$ with $\pi'.\mathrm{length} \ge 1$. By (E), $\pi'^{\uparrow} : v \tuh \cdots \tuh v$ is a directed walk in $G$ of length $\ge \pi'.\mathrm{length} \ge 1$ (each shortcut expands to a sub-walk of $G$ of length $\ge 1$, and lengths add). Moreover, $v \in G^{\sm W}$ gives $v \in J \cup (V \sm W) \ins J \cup V = G$, so $v \in G$ and there is a non-trivial directed cycle through $v$ in $G$, contradicting \refrow{def_3_6}.

\subsubsection*{Item 3b -- topological order.}
\textit{(Justifies \texttt{marginalize\_isTopologicalOrder}.)} Let $r$ be a topological order of $G$ (\refrow{def_3_8}); we show $r$ is also a topological order of $G^{\sm W}$ -- the LN's ``by just ignoring the nodes from $W$'' is precisely the statement that the \emph{same} relation $r$, restricted to the smaller vertex set of $G^{\sm W}$, satisfies the topological-order axioms.

Every $v \in G^{\sm W}$ satisfies $v \in J \cup (V \sm W) \ins J \cup V = G$. The first three axioms of \refrow{def_3_8} (irreflexivity, transitivity, trichotomy) on $G^{\sm W}$ then follow by restriction from the corresponding axioms on $G$.

For parent precedence, let $v, w \in G^{\sm W}$ with $v \in \Pa^{G^{\sm W}}(w)$, i.e.\ $(v, w) \in E^{\sm W}$. The $E^{\sm W}$ characterisation supplies a directed walk
\[
   v = t_0 \tuh t_1 \tuh \cdots \tuh t_{n-1} \tuh t_n = w \qquad\text{in } G
\]
of length $n \ge 1$ with $t_1, \dots, t_{n-1} \in W$. Each step $(t_i, t_{i+1}) \in E$ gives $t_i \in \Pa^G(t_{i+1})$, so $r\,t_i\,t_{i+1}$ by the parent-precedence axiom of $r$ on $G$ (each $t_i \in G$: the endpoints lie in $G^{\sm W} \ins G$, the intermediates lie in $V \ins G$). Chaining through $r$'s transitivity on $G$ yields $r\,v\,w$.
\end{proof}

\end{document}
